\chapter{Build, buy or combine existing methods?}
\label{exec16:s10}

Chapter~\ref{scisyn:software} gives the principal software recommendation.
This chapter supplies its comparative decision argument: admission before
price, paired outcomes on the assigned population, bounded qualification and
maintenance-inclusive cost. The final DataRobot study applies those criteria
to a concrete propensity-modeling workflow, distinguishing a useful component
from a complete causal estimator.

\section{Choose the method before the supplier}

A bank selecting causal software is choosing a procedure for a specified
decision, together with the means of maintaining and operating that procedure.
Discovery for a supplied independent sample, identification in a latent graph,
and evaluation of a repeated servicing policy are different services. The
software survey establishes candidate routes and their restrictions. This
chapter asks how those routes should be compared once a bounded use has been
chosen, and what evidence could justify direct adoption, a process bridge, a
port, an internal implementation, a commercial purchase, or reference-only use.

Let \(q\) denote the complete requested scientific operation and \(s\) a
declared workload scenario. For candidate route \(j\), write
\[
 a_j(q,s)=a_j^{\rm method}(q,s)\,a_j^{\rm representation}(q,s)
          \,a_j^{\rm operating}(q,s)\,a_j^{\rm rights}(q,s),
\]
where each factor is one only when the corresponding necessary condition has
been established in the declared scope. The factors describe compatibility
with the causal method, preservation of its objects, the required runtime
behavior, and the intended reuse arrangement. An unknown condition remains
unresolved; it is not a passing value. This is an admission rule for a proposed
comparison, not a probability model or a weighted capability score. A low
price cannot compensate for a method that answers a different intervention
query, and a favorable license cannot supply missing uncertainty.
\citep{exec16_s10_design}

Admission should distinguish hard scientific restrictions from tunable
engineering requirements. If the selected routine only handles independent
units, declaring a panel identifier does not make it a dependent-data method.
If the issue is process startup, an alternative process lifecycle may resolve
it without changing the target. If the issue is an unsupported counterfactual,
an adapter cannot resolve it by relabelling an associational prediction. This
classification tells the proposal whether to restrict the supported domain,
implement a conversion, select another scientific method, or commission
research. It prevents an implementation task from carrying an implicit promise
to solve nonidentification.

The inspected candidates illustrate why a route is more specific than a
package name. The pinned pgmpy review finds an advertised adjustment option
whose selected implementation raises \texttt{NotImplementedError}; that option
cannot enter the executable comparison merely because it appears in a guide.
The selected Tetrad review supports a restricted Java PC prototype subject to
runtime, conversion and reuse conditions, not an unrestricted Java discovery
service. The CausalNex review treats the inspected discontinued source as an
archival comparison; maintenance and unreported solver behavior remain part
of the decision. These are findings for inspected revisions. They do not
establish that another release, method or supplier has the same restriction.
\citep{exec13_pgmpy_reuse_primary,exec13_tetrad_reuse_decision,exec14_causalnex_primary}

\section{Matching the experiment and its denominator}

A comparison requires a common request population before it requires a common
timer. Let \(\mathcal Q\) be the prespecified distribution of requested cases,
including difficult and unsupported cases relevant to the proposed service.
For each route, retain a disposition \(R_j(q)\) such as computed, unsupported,
unavailable, or invalid input, together with the target and uncertainty object
actually returned. The intended denominator is the requested workload. A
conditional accuracy figure among successful outputs remains useful, but it
answers a different question from the probability of receiving a correct,
usable result for an arbitrary requested case.

For estimator comparisons, a matched case must bind the population measure,
intervention, horizon, information clock and uncertainty meaning. The two
procedures may legitimately use different computational strategies. They may
also use different statistical methods if the purpose is a method comparison,
provided the target and admissible information are held fixed and each method's
assumptions are exposed. A port-equivalence study has a narrower objective: it
tests whether the intended procedure survives translation, including defaults,
failure behavior, randomness and diagnostic outputs. Mixing these objectives
would make a better result ambiguous between a scientific change and a faithful
implementation.

The first prototype should contain three complementary groups. Exact finite
or analytic constructions test meanings and numerical identities. Generated
sampling experiments test operating characteristics under explicit truth.
Documented public datasets test preparation, interfaces, and useful questions
within their evidence limits. Agreement on public observational data is not
itself a causal-truth label. A benchmark with known intervention assignments
supports a stronger comparison only for its documented population and design.
The proposal can therefore establish the merit and practicality of a later
bank experiment without pretending that public data reproduces private bank
assignment, identity or operational truth.
\citep{exec16_s10_design}

Workload matching also includes resource and information budgets. Both routes
should receive the same permitted data and expert restrictions; preprocessing,
search, tuning and human repair must be charged consistently. If one candidate
uses a larger graph-search budget, report both a fixed-budget comparison and
the declared cost--quality tradeoff. For a bridge, include serialization, data
movement, environment startup, numerical work, uncertainty construction and
publication. Warm and cold executions are distinct operating scenarios. Reuse
of a fitted model or cached result is valid only when the identity and
invalidation rule agree with the intended service.

\section{A paired comparison can reverse a preferred route}

Consider an authored workload law with 1,000 equally weighted request slots.
Each route produces a correct usable answer, a wrong answer, or an unavailable
answer. The following table specifies the entire paired law; it is not an
observed package trial. A row refers to route \(A\), and a column to route
\(B\).
\[
\begin{array}{c|rrr|r}
 &B:\mathrm{correct}&B:\mathrm{wrong}&B:\mathrm{unavailable}&\mathrm{total}\\\hline
 A:\mathrm{correct}&600&20&10&630\\
 A:\mathrm{wrong}&10&5&5&20\\
 A:\mathrm{unavailable}&190&75&85&350\\\hline
 \mathrm{total}&800&100&100&1000
\end{array}
\]
Among available answers, \(A\) is correct in \(630/650=126/130\) cases and
\(B\) in \(800/900=8/9\). On the requested denominator, their correct-answer
rates are 0.63 and 0.80. Thus \(A\) has higher conditional correctness while
\(B\) supplies more correct answers overall. Neither statement alone selects
the preferred banking workflow.
\citep{exec16_s10_exact}

Suppose a correct answer has stipulated gross value USD 10, a wrong answer
causes loss \(\ell\), unavailable answers have zero gross value, and the
complete per-request costs are USD 1 for \(A\) and USD 2 for \(B\), charged on
all requests. Their total values are
\[
 V_A=630(10)-20\ell-1000,\qquad
 V_B=800(10)-100\ell-2000.
\]
The difference is \(V_B-V_A=700-80\ell\). Route \(B\) is preferred only when
\(\ell<8.75\). At \(\ell=40\), the totals are USD 4,500 and USD 2,000,
respectively. This threshold is a sensitivity result under a specified loss
and fallback convention. It is not a market price or a claim about a bank's
actual error cost. If unavailable answers require manual review, that fallback
has its own value, delay and cost and must be incorporated in both workflows.
\citep{exec16_s10_exact}

Pairing provides information that two separate success percentages lose. Let
\(X_j\) be the indicator of a correct usable answer from route \(j\), and let
\(D=X_B-X_A\). In the stipulated law, \(D=1\) on 200 slots, \(D=-1\) on 30,
and \(D=0\) on 770. Hence
\[
 E(D)=0.17,\qquad
 \operatorname{Var}(D)=0.23-0.17^2=0.2011.
\]
Equivalently, the covariance of the two correctness indicators is
\(0.600-(0.63)(0.80)=0.096\), and
\[
 \operatorname{Var}(D)=0.63(0.37)+0.80(0.20)-2(0.096).
\]
If independent paired requests were drawn from this law, the sample mean
difference would have variance \(0.2011/n\). Treating the two outcomes as
independent instead gives \(0.3931/n\). These expressions describe the
stipulated sampling experiment; they are not confidence statements for a
fixed benchmark table. Repeated variants from one customer, graph or source
document require their actual group dependence to be retained.
\citep{exec16_s10_exact}

The paired table also makes disagreement review efficient. The 200 cases on
which only \(B\) is correct and the 30 on which only \(A\) is correct are
natural targets for diagnosing method domain, conversion or failure handling.
They should remain in the final outcome accounting after investigation.
Correcting a test oracle or an adapter creates a successor comparison with
explicit changes; removing disagreeing cases until the routes agree would
change the workload and obscure the reason for the improvement.

\section{What a failure-free prototype establishes}

A set of successful fixtures proves the stated examples, not universal
compatibility. A small stochastic qualification experiment supplies a
different, conditional statement. Suppose a fixed workload distribution has
independent Bernoulli failures with probability \(p\), and the evaluation
contains \(n\) prespecified requests. If zero failures occur, the probability
of that observation is \((1-p)^n\). Inverting the one-sided test at error
probability \(\alpha\) gives the upper bound
\[
 p_U=1-\alpha^{1/n}.
\]
With \(n=100\) and \(\alpha=0.05\), this is approximately 0.0295. Zero failures
among a hundred independent requests therefore does not establish a one-percent
failure ceiling at this level. The exact inequality
\[
 (0.99)^n\leq0.05
\]
first holds at \(n=299\). This is a prospective sample-size implication if
that design and ceiling are appropriate, not a claim that the prototype has
passed it.
\citep{exec16_s10_exact}

For eight separately declared profiles, a simultaneous statement needs its
own error allocation. Assigning \(\alpha/8\) to each and applying the union
bound yields joint error at most \(\alpha\) without requiring independence
between profiles. Each within-profile zero-failure calculation still requires
its stated independent Bernoulli experiment. At a common one-percent ceiling,
\((0.99)^n\leq0.05/8\) first holds at \(n=505\) requests per profile. A single
pooled bound would answer a different question and could hide a bad rare
profile. Conversely, infinitely many graph families or deployment regimes
cannot be qualified by calling eight chosen profiles exhaustive.
\citep{exec16_s10_exact}

These calculations motivate a bounded qualification protocol rather than a
large unstructured test collection. Define the admitted profile, failure event,
request distribution, denominator, sample size, simultaneous statement and
stopping rule before evaluation. A crash, timeout, missing interval or wrong
target is not automatically the same event; retain their separate frequencies
and any combined service failure used for acceptance. If a generated experiment
tests confidence-set coverage, its event is whether the known fixed target is
contained, together with declared availability and precision requirements.
Statistical coverage, computation completion and useful decision resolution
need separate criteria. Adaptive test selection or optional stopping requires
a design that justifies it; repeatedly adding cases until the displayed bound
passes does not preserve this fixed-sample argument.

An exact finite-domain test can offer stronger evidence within a smaller
domain. Exhausting all supported endpoint pairs or all three-node graphs
may settle the declared finite conversion behavior. It does not become a
sampling confidence claim, nor does it establish arbitrary-dimensional
algorithm correctness. The topology examples show how exact counterexamples
and actual codecs complement rather than replace a generated performance
study. The proposal should state which of these evidence forms closes each
acceptance condition.
\citep{exec16_topology_exact,exec16_topology_codec}

\section{A port must preserve the full procedure}

Porting has several possible meanings. A mathematical reimplementation may
reproduce a functional using a different algorithm. A compatibility port may
seek the same defaults, accepted inputs, warnings and output semantics as a
specified release. A new method may deliberately improve an estimator or
uncertainty construction. Each can be useful, but each requires a different
comparison. The product should not retain an upstream method name while
changing its weighting measure, policy target or interval construction without
declaring that change. Source-level agreement is particularly weak when a
caller changes the meaning of an otherwise correctly translated routine.

For the inspected CausalPy route, the proposed sharp-discontinuity contract
distinguishes a finite-epsilon prediction contrast from a threshold estimand,
and a point backend from posterior draws and predictive responses. A port that
reproduces a plotted difference while discarding epsilon, treatment direction
or left/right support has not reproduced the causal operation. For the d3rlpy
route, the proposed contract separates episode termination, administrative
cutoff, missing continuation, initial-state population, policy identity and
the physical discount clock. Reproducing an optimizer while averaging the
wrong initial states can produce a different policy value. These source-grounded
contracts specify tests to execute before enabling the proposed integrations;
they are not records of native acceptance trials.
\citep{exec14_causalpy_decision,exec15_d3rlpy_contract}

Numerical acceptance should be stated in the coordinates of the decision.
Let \(\widetilde\theta\) approximate an independently defined reference
\(\theta\), with justified deterministic bounds
\(\lvert\widetilde\theta_i-\theta_i\rvert\leq\varepsilon_i\). For fixed
weights \(w_i\) and cost \(c\), the triangle inequality gives
\[
 \left|\sum_iw_i\widetilde\theta_i-c-
              \left(\sum_iw_i\theta_i-c\right)\right|
 \leq\sum_i|w_i|\varepsilon_i.
\]
This transports certified numerical error; it does not certify the reference's
causal validity or its sampling uncertainty. Relative-error tolerances alone
are unsuitable near a zero-valued target, and coordinate tolerances that look
small can still reverse a near-threshold decision.

For a stipulated example, let \(w=(3/4,1/4)\),
\(\widetilde\theta=(24,-8)\) USD, \(c=15\) USD and
\(\varepsilon=(2,4)\) USD. The approximate net value is USD 1, but its
numerical error allowance is USD 2.5, giving the enclosure
\([-1.5,3.5]\). The permitted reference \(\theta=(23,-10)\) yields USD
\(-0.25\). Declaring a profitable action from the approximate point would
exceed the available numerical evidence. Reducing the numerical budget,
retaining an unresolved decision, or changing the information/design are
distinct responses. If sampling uncertainty is also present, the numerical
enclosure must be combined with that construction under its actual guarantee,
not relabelled as another confidence level.
\citep{exec16_s10_exact}

Port acceptance should consequently include independent algebraic references,
metamorphic transformations with known consequences, adverse inputs, and
method-specific numerical error checks. Graph and set outputs need structural
and containment criteria, not only a mean squared difference. Seeds, thread
counts, solver versions, hardware-dependent arithmetic and stopping conditions
belong in the tested procedure. An upstream test suite is valuable inherited
evidence, but common implementation mistakes can survive a direct translation
of both code and tests. Independent references should target the scientific
properties the proposed product promises.

\section{Lifecycle cost conditional on qualified behavior}

Only admitted routes should enter a procurement comparison. For route \(j\),
a transparent planning quantity over \(T\) periods is
\[
 C_j=F_j+\sum_{t=1}^{T}d_t
  \{L_{jt}+M_{jt}+N_t c_{jt}+I_{jt}\},
\]
where \(F_j\) is initial engineering and qualification cost, \(L_{jt}\)
license or subscription expense, \(M_{jt}\) maintenance and compatibility
work, \(N_t c_{jt}\) workload-dependent cost, and \(I_{jt}\) the expected
incident or recovery cost under an explicitly stated model. The \(d_t\)
are the chosen discount factors. These terms need not be known at proposal
time. Ranges and break-even conditions are more informative than an invented
precise total, and omitted costs must remain visible.

Suppose, purely for planning, a bridge costs USD 100,000 initially, USD 20,000
per year to maintain and USD 3 per request. A port costs USD 300,000 initially,
USD 80,000 per year and USD 1 per request. Ignore discounting over three years
and assume both routes meet exactly the same service requirements. Ignoring
maintenance suggests a break-even volume of
\[
 \frac{300000-100000}{3-1}=100000
\]
requests. Including the three years of maintenance gives
\[
 \frac{300000-100000+3(80000-20000)}{3-1}=190000.
\]
At 150,000 total requests the bridge costs USD 610,000 and the port USD
690,000. The cheaper numerical kernel does not establish the cheaper full
route. These are exact hypothetical calculations, not measured costs, vendor
quotes or return on investment. Their practical use is to identify which
workload and maintenance estimates could change the decision.
\citep{exec16_s10_exact}

The same decomposition organizes substantive alternatives. Direct adoption
can minimize translation work when the upstream environment and result contract
fit, but it retains dependency and upgrade obligations. A process bridge can
isolate an environment and preserve the upstream implementation, at the cost
of conversion, lifecycle and operational work. A port may serve an unavailable
deployment environment or a demonstrated bottleneck, but incurs full-procedure
equivalence and maintenance responsibilities. An internal implementation is
justified when the required scientific construction or result type is absent
from suitable existing routes; its comparison should include the continuing
cost of independent verification. Reference-only use is appropriate when an
implementation clarifies a method but its runtime, rights, maintenance or
scientific domain is unsuitable for the proposed service.

A commercial purchase is evaluated through the same admitted operation and
whole-workflow cost. The exact edition, automation route, technical method,
diagnostics, export semantics and intended deployment terms matter. Community
extensions should not be attributed to the vendor's supported core without
evidence. Missing technical documentation leaves a capability unresolved; it
does not establish absence. Unquoted prices remain unknown. The commercial
comparison therefore supports a request for a matched evaluation or a bounded
purchase decision when its evidence permits one, rather than an unsupported
claim of product superiority.
\citep{exec16_s10_design}

\section{A proposal decision that can be executed}

A provisional choice should specify the method, environment and comparison that could justify adoption. A useful
decision identifies the candidate revision and environment, the exact causal
operation, the alternative route, the information and compute budget, the
correctness and failure events, the required numerical or uncertainty
guarantee, and the cost quantities that remain to be measured. It also states
what outcome would change the choice. A preferred prototype candidate is
therefore compatible with deferred integration; it expresses where the next
evidence is most valuable without asserting that the evidence exists.
\citep{exec16_s10_design}

For a discovery prototype, a restricted supplied regime and independently
checked graph conversions precede an effect claim. For a threshold design,
the treatment direction, support and requested limit precede a comparison of
predictions. For servicing, the complete policy and episode population precede
offline value comparison. For funding, weak information and the topology of
the effect set precede a monetary recommendation. These conditions let
independent prototype routes advance while another method remains unsupported.
They also make rejection useful: a documented mismatch removes a route from
the admitted domain and identifies whether new information, another method or
engineering work could restore it.

No proposed acceptance test can make observationally indistinguishable causal
models distinguishable by computation alone. Nor can it invent bank precision
targets, authentic identity and assignment records, supplier commitments or
independent approval. Those limits should govern the claim made by the
proposal, while public/generated prototype comparisons establish what can be
learned now. A later bank trial is justified by a credible scoped method,
demonstrated semantic and numerical behavior, and a decision-relevant
uncertainty and cost question. The resulting build, buy or integrate decision
is conditional on that evidence, with an explicit fallback when qualification
fails. The following propensity-modeling study applies that rule to a
specific platform workflow: the prediction component is useful only if its
training information, row identity, weighting target and uncertainty survive
the proposed integration.
\label{exec16:s10-end}
